\section{Method}
\label{sec:method}

\subsection{Task Formulation and Auditable Data Protocol}
We consider visually observable manipulation in which the object is already
grasped, the target or task state is visible, and successful execution requires
local trajectory alignment or refinement. An episode is
\begin{equation}
  \tau=\{(o_t,u^{\mathrm{raw}}_t,u^{\mathrm{exec}}_t,q_t,e_t)\}_{t=1}^{T},
  \label{eq:episode}
\end{equation}
where $o_t$ denotes synchronized multi-camera observations,
$u^{\mathrm{raw}}_t$ the operator command, $u^{\mathrm{exec}}_t$ the command
applied by the robot controller, $q_t$ the measured robot state, and $e_t$ a
timestamped audit event. Each task configuration declares its success
criterion, pose tolerance, velocity limit, and safety termination condition
before collection. Grasp acquisition, long-horizon navigation, force-dependent
control without force sensing, persistent occlusion, and deformable-object
dynamics are outside our present scope.

An annotator marks one or more correction intervals
$I_k=[t_s^k,t_e^k]$, spanning the onset of a corrective maneuver through its
recovery. Their union induces the frame-level indicator
\begin{equation}
  m_t=\mathbb{1}\!\left[t\in\bigcup_k I_k\right]\in\{0,1\}.
  \label{eq:correction_indicator}
\end{equation}
This indicator is derived from interval annotation and is used only to build
offline supervision and evaluation statistics. The episode contract also
records terminal success or failure, manual bypass, safety termination, and
task-configuration metadata.

Let $\Gamma(\tau)$ be a fixed admission rule over terminal outcome,
synchronization completeness, missing frames, and safety events. Successful,
reusable episodes satisfying this rule form the action-training set
\begin{equation}
  \Aaction=\{\tau:\Gamma(\tau)=1\},
\end{equation}
whereas failed, near-miss, conflicting, incomplete, or safety-terminated
episodes are retained in $\Aaudit$ for quality analysis. Correction intervals
inside admitted episodes supply difficult-state behavior; their successful
nominal complement $\mathcal{N}=\{t:m_t=0\}$ identifies contexts in which the
operator command should be preserved. The processing path is
\emph{raw record $\rightarrow$ synchronized episode $\rightarrow$ audited
training view $\rightarrow$ collection model}.

\begin{figure}[t]
  \centering
  \resizebox{\columnwidth}{!}{%
  \begin{tikzpicture}[font=\sffamily\scriptsize,>=Latex]
    \def\xmax{10.8}

    % Correction interval background shared by all panels.
    \fill[GainOrange!13] (3.25,-0.3) rectangle (7.15,5.55);
    \draw[GainOrange!70,dashed] (3.25,-0.3) -- (3.25,5.55);
    \draw[GainOrange!70,dashed] (7.15,-0.3) -- (7.15,5.55);
    \node[text=GainOrange!90!black,font=\sffamily\scriptsize\bfseries]
      at (5.2,5.3) {correction interval $I_k$};

    % Interval label.
    \draw[->,draw=DarkGray] (0,4.35) -- (\xmax,4.35);
    \draw[->,draw=DarkGray] (0,4.1) -- (0,5.05);
    \draw[line width=1.3pt,draw=AuditBlue]
      (0.1,4.45) -- (3.25,4.45) -- (3.25,4.88) --
      (7.15,4.88) -- (7.15,4.45) -- (10.55,4.45);
    \node[anchor=east] at (-0.15,4.7) {$m_t$};
    \node[anchor=north] at (1.65,4.38) {nominal};
    \node[anchor=north] at (8.85,4.38) {nominal};

    % Continuous gain.
    \draw[->,draw=DarkGray] (0,2.45) -- (\xmax,2.45);
    \draw[->,draw=DarkGray] (0,2.2) -- (0,3.75);
    \draw[line width=1.5pt,draw=GainOrange]
      plot[smooth] coordinates {(0.1,2.55) (1.4,2.57) (2.7,2.63)
      (3.5,2.82) (4.3,3.25) (5.2,3.55) (6.2,3.48)
      (7.1,3.12) (8.0,2.75) (9.2,2.58) (10.55,2.56)};
    \node[anchor=east,text=GainOrange!90!black] at (-0.15,3.1) {$\alpha_t$};
    \node[anchor=west,text=GainOrange!90!black] at (7.45,3.35)
      {bounded rate};

    % Raw and assisted action schematic.
    \draw[->,draw=DarkGray] (0,0.25) -- (\xmax,0.25);
    \draw[->,draw=DarkGray] (0,-0.05) -- (0,1.75);
    \draw[line width=1.1pt,draw=DarkGray!65]
      plot[smooth] coordinates {(0.1,0.62) (1.2,0.9) (2.2,0.55)
      (3.2,0.85) (4.1,0.38) (5.0,1.35) (5.7,0.48)
      (6.5,1.18) (7.3,0.75) (8.4,1.02) (9.4,0.72) (10.55,0.9)};
    \draw[line width=1.5pt,draw=ActionTeal]
      plot[smooth] coordinates {(0.1,0.65) (1.2,0.84) (2.2,0.62)
      (3.2,0.78) (4.1,0.68) (5.0,0.92) (5.7,0.73)
      (6.5,0.98) (7.3,0.82) (8.4,0.97) (9.4,0.77) (10.55,0.88)};
    \node[anchor=east] at (-0.15,1.05) {action};
    \draw[line width=1.1pt,draw=DarkGray!65] (7.65,1.55) -- (8.25,1.55);
    \node[anchor=west] at (8.35,1.55) {raw};
    \draw[line width=1.5pt,draw=ActionTeal] (7.65,1.25) -- (8.25,1.25);
    \node[anchor=west] at (8.35,1.25) {assisted};

    \node[anchor=north] at (\xmax,0.18) {time};
  \end{tikzpicture}}
  \caption{Schematic relation between interval supervision and continuous
  assistance. Binary interval annotations define an offline training signal;
  the deployed gain remains a bounded, rate-constrained continuous state. The
  action trace is illustrative and does not report experimental data.}
  \label{fig:supervision_gain}
\end{figure}


\subsection{Visual-Conditioned Continuous Action Filter}
At time $t$, the filter uses a causal history of length $L$,
\begin{equation}
  x_t=\{u^{\mathrm{raw}}_{t-L:t-1},q_{t-L:t-1},v_{t-L:t}\},
  \qquad v_t=E_{\mathrm{vis}}(o_t),
  \label{eq:context}
\end{equation}
where the frozen encoder $E_{\mathrm{vis}}$ maps the synchronized camera views
to a continuous representation $v_t$. We use SigLIP2 for this representation.
An offline vision-language auditor processes longer clips and proposes task
phase, progress, stagnation, recovery, and correction events with confidence
scores. Human review converts these proposals into interval labels, terminal
labels, and the admission decision. Online inference uses only the context in
Eq.~\eqref{eq:context}.

Linear projections convert command, state, and visual features into tokens. A
causal Transformer maps their interleaved history to $h_t$ using the attention
mask $M_{ij}=0$ for $j\leq i$ and $-\infty$ otherwise
\citep{vaswaniAttentionAllYou2023}. The action branch models a conditional
expert-action distribution with a conditional variational autoencoder (CVAE)
\citep{kingmaAutoEncodingVariationalBayes2022}:
\begin{align}
  z_t &\sim p_\theta(z_t\mid h_t), \\
  \hat u^{\mathrm{exp}}_t &\sim
  p_\theta(u^{\mathrm{exp}}_t\mid h_t,z_t),
  \label{eq:cvae_action}
\end{align}
with an amortized posterior
$q_\phi(z_t\mid h_t,u^{\mathrm{exp}}_t)$ used during training. The latent
variable represents multiple locally plausible expert actions under similar
visual context.

A separate gain branch maps $h_t$ to a bounded increment and updates a scalar
assistance state,
\begin{align}
  \Delta\alpha_t &= r_\alpha\tanh(g_\theta(h_t)), \\
  \alpha_t &= \operatorname{clip}
  (\alpha_{t-1}+\Delta\alpha_t,0,\alpha_{\max}),
  \label{eq:gain}
\end{align}
where $r_\alpha$ limits the per-step gain change and $\alpha_{\max}$ limits
authority. The candidate expert action defines the local residual and the gain
sets its applied magnitude:
\begin{align}
  \hat\delta_t &= \hat u^{\mathrm{exp}}_t-u^{\mathrm{raw}}_t, \\
  u^{\mathrm{out}}_t &= \Pi_{\mathrm{safe}}\!\left(
  u^{\mathrm{raw}}_t+\alpha_t\hat\delta_t\right).
  \label{eq:composition}
\end{align}
$\Pi_{\mathrm{safe}}$ projects the composed command onto declared joint
position, velocity, and command-rate constraints. The executed action changes
$o_{t+1}$, so Eqs.~\eqref{eq:context}--\eqref{eq:composition} recur on the
resulting visual history.

\subsection{Complementary Supervision}
For synchronized trajectories in $\Aaction$, let
$u^{\mathrm{exp}}_t$ denote the recorded expert command. We optimize
\begin{equation}
  \mathcal{L}=\lambda_a\mathcal{L}_{\mathrm{act}}+
  \lambda_g\mathcal{L}_{\mathrm{gain}}+
  \lambda_n\mathcal{L}_{\mathrm{nom}}+
  \beta\mathcal{L}_{\mathrm{KL}}+
  \lambda_s\mathcal{L}_{\mathrm{smooth}}.
  \label{eq:objective}
\end{equation}
The action term emphasizes expert behavior inside correction intervals,
\begin{align}
  \mathcal{L}_{\mathrm{act}}
  &=\frac{1}{T}\sum_t w_t
  \|\hat u^{\mathrm{exp}}_t-u^{\mathrm{exp}}_t\|_1,\\
  w_t&=1+(w_{\mathrm{corr}}-1)m_t.
  \label{eq:action_loss}
\end{align}
The gain term supplies interval-level trend supervision rather than a target
gain at every frame:
\begin{equation}
\begin{split}
  \mathcal{L}_{\mathrm{gain}}=\frac{1}{T}\sum_t
  \big[&m_t[\alpha_{\min}-\alpha_t]_+\\
       &+(1-m_t)\alpha_t\big],
  \label{eq:gain_loss}
\end{split}
\end{equation}
where $[a]_+=\max(a,0)$ and $\alpha_{\min}$ denotes the minimum useful gain
range selected on validation data. Successful nominal segments further
penalize candidate corrections that would move an already adequate command,
\begin{equation}
  \mathcal{L}_{\mathrm{nom}}=\frac{1}{T}\sum_t(1-m_t)
  \|\hat\delta_t\|_1.
  \label{eq:nominal_loss}
\end{equation}
The CVAE regularizer is
\begin{equation}
  \mathcal{L}_{\mathrm{KL}}=
  D_{\mathrm{KL}}\!\left(q_\phi(z_t\mid h_t,u^{\mathrm{exp}}_t),
  \|,p_\theta(z_t\mid h_t)\right),
\end{equation}
and $\mathcal{L}_{\mathrm{smooth}}$ penalizes adjacent changes in both
$u^{\mathrm{out}}_t$ and $\alpha_t$. This factorization lets the action branch
represent where a plausible local correction leads while the gain branch
controls how much of it enters the command stream.

\subsection{Safe Iterative Data Collection}
Online collection executes the encoder, temporal model, action and gain
branches, and safety projection. Runtime checks cover non-finite outputs,
residual magnitude and rate, joint and velocity limits, model timeout, and
missing visual frames; the operator retains a manual bypass and hardware
emergency stop.

Before round zero, we freeze a reference audit set $H_{\mathrm{ref}}$ and the
admission rule $\Gamma$. At round $r$, dataset $D_r$ trains filter $f_r$, which
collects candidate episodes $D'_r$. A configuration-stratified subset
$H^{(r)}_{\mathrm{audit}}$ is independently re-audited and permanently excluded
from training. The remaining admitted candidates update the next round:
\begin{equation}
\begin{split}
  D_r &\xrightarrow{\mathrm{train}} f_r
  \xrightarrow{\mathrm{collection}} D'_r,\\
  D_{r+1} &=D_r\cup
  \Gamma\!\left(D'_r\setminus H^{(r)}_{\mathrm{audit}}\right).
  \label{eq:flywheel}
\end{split}
\end{equation}
$H_{\mathrm{ref}}$ detects cross-round audit drift, while
$H^{(r)}_{\mathrm{audit}}$ estimates false admission and missed valid episodes
on newly collected data. With $\Gamma$, operator time, reset count, task quotas,
and safety limits fixed, changes in admitted demonstrations per operator minute
measure the practical effect of each update.
